\section*{Влияние волатильности на цену опциона}

\subsection*{Цель работы}

Исследовать влияние волатильности на цену европейского call-опциона на примере акций 20 компаний с использованием формулы Блэка--Шоулза.

\subsection*{Теоретические сведения}

Цена европейского call-опциона по модели Блэка--Шоулза определяется формулой:

\[
C = S_0 N(d_1) - Ke^{-rT}N(d_2),
\]

где

\[
d_1 = \frac{\ln\left(\frac{S_0}{K}\right)+\left(r+\frac{\sigma^2}{2}\right)T}{\sigma\sqrt{T}},
\]

\[
d_2 = d_1 - \sigma\sqrt{T}.
\]

Здесь:
\begin{itemize}
    \item \(S_0\) --- текущая цена акции;
    \item \(K\) --- цена исполнения опциона;
    \item \(r\) --- безрисковая процентная ставка;
    \item \(T\) --- время до исполнения;
    \item \(\sigma\) --- волатильность акции;
    \item \(N(\cdot)\) --- функция стандартного нормального распределения.
\end{itemize}

\subsection*{Исходные параметры}

Для всех компаний примем:
\[
K=S_0,
\]
то есть опцион находится ``at the money''.

Также зададим:
\[
r=5\%, \qquad T=1 \text{ год}.
\]

\subsection*{Данные по компаниям}

\begin{center}
\begin{tabular}{|c|c|c|}
\hline
№ & Компания & Волатильность \(\sigma\) \\
\hline
1 & Apple & 0.20 \\
2 & Microsoft & 0.18 \\
3 & Amazon & 0.30 \\
4 & Google & 0.22 \\
5 & Tesla & 0.55 \\
6 & NVIDIA & 0.45 \\
7 & Meta & 0.35 \\
8 & Intel & 0.28 \\
9 & Netflix & 0.40 \\
10 & Coca-Cola & 0.15 \\
11 & PepsiCo & 0.16 \\
12 & JPMorgan & 0.25 \\
13 & Goldman Sachs & 0.27 \\
14 & Boeing & 0.38 \\
15 & Nike & 0.24 \\
16 & McDonald's & 0.17 \\
17 & Walmart & 0.14 \\
18 & Disney & 0.29 \\
19 & IBM & 0.19 \\
20 & Adobe & 0.31 \\
\hline
\end{tabular}
\end{center}

\subsection*{Расчёт цен опционов}

Пусть для всех компаний:
\[
S_0 = K = 100.
\]

Тогда вычислим цену call-опциона для различных значений волатильности.

\begin{center}
\begin{tabular}{|c|c|c|}
\hline
Компания & Волатильность \(\sigma\) & Цена опциона \(C\) \\
\hline
Apple & 0.20 & 10.45 \\
Microsoft & 0.18 & 9.68 \\
Amazon & 0.30 & 14.23 \\
Google & 0.22 & 11.20 \\
Tesla & 0.55 & 24.87 \\
NVIDIA & 0.45 & 20.65 \\
Meta & 0.35 & 16.48 \\
Intel & 0.28 & 13.51 \\
Netflix & 0.40 & 18.02 \\
Coca-Cola & 0.15 & 8.59 \\
PepsiCo & 0.16 & 8.96 \\
JPMorgan & 0.25 & 12.34 \\
Goldman Sachs & 0.27 & 13.10 \\
Boeing & 0.38 & 17.32 \\
Nike & 0.24 & 11.91 \\
McDonald's & 0.17 & 9.31 \\
Walmart & 0.14 & 8.15 \\
Disney & 0.29 & 13.87 \\
IBM & 0.19 & 10.06 \\
Adobe & 0.31 & 14.76 \\
\hline
\end{tabular}
\end{center}

\subsection*{Анализ результатов}

Из полученных данных видно, что с увеличением волатильности возрастает стоимость опциона.

Например:
\begin{itemize}
    \item при волатильности \(0.14\) цена опциона составляет около \(8.15\);
    \item при волатильности \(0.30\) цена возрастает до \(14.23\);
    \item при волатильности \(0.55\) цена достигает \(24.87\).
\end{itemize}

Это связано с тем, что высокая волатильность увеличивает вероятность значительного роста цены базового актива, тогда как убытки покупателя опциона ограничены уплаченной премией.

\subsection*{Вывод}

В ходе работы было исследовано влияние волатильности на цену европейского call-опциона.

На примере 20 компаний показано, что:
\begin{itemize}
    \item цена опциона положительно зависит от волатильности;
    \item чем выше риск и изменчивость акции, тем дороже опцион;
    \item волатильность является одним из ключевых параметров модели Блэка--Шоулза.
\end{itemize}

Следовательно, рост волатильности приводит к увеличению стоимости опциона.
